% In this file you should put the actual content of the blueprint.
% It will be used both by the web and the print version.
% It should *not* include the \begin{document}
%
% If you want to split the blueprint content into several files then
% the current file can be a simple sequence of \input. Otherwise It
% can start with a \section or \chapter for instance.

\chapter{Introduction}

note: this blueprint was drafted by Claude

\emph{Moser's worm problem} asks for the (convex) region of smallest area in
the plane that contains a rotated and/or translated copy of every curve of unit
length (a \emph{worm}). Any such region is called a \emph{Moser set}. The
exact optimum is unknown; the best published upper and lower bounds sandwich
it between $0.232239$ and $\pi/12$. 

This project pursues a \emph{computer-checked} improvement to the lower bound, by showing that every Moser set must contain a convex polygon of area at least $0.232240$. 

The main idea is to discretise the problem. We create and iteratively mutate a \emph{working set} of convex
candidate polygons satisfying the property that which any convex Moser set of size below $0.232240$ must contain a copy of at least one of the candidates. 
The working set is initialised with a single polygon, and then repeatedly mutated
by replacing elements with collections of larger elements in ways that preserve the property, all the while removing any element of size greater than $0.232240$.
If we can find a sequence of mutations that ends with an empty working set, we may conclude that no convex Moser set of area below $0.232240$ exists.

\chapter{Rational planar geometry}

All geometry in this project is carried out over $\mathbb{Q}^2$ so that
decidable computations can be used throughout. Real-valued errors enter
only through explicit approximation bounds.

\begin{definition}[Rational point]
  \label{def:rational-point}
  \lean{RationalPoint}
  \leanok
  A \emph{rational point} is an element of $\mathbb{Q}^2$, represented as a
  function $\mathrm{Fin}\,2 \to \mathbb{Q}$.
\end{definition}

\begin{definition}[Cross product, dot product, squared length]
  \label{def:rational-ops}
  \lean{RationalPoint.crossProduct, RationalPoint.dotProduct, RationalPoint.lengthSq}
  \leanok
  \uses{def:rational-point}
  For $u, v \in \mathbb{Q}^2$ we define
  $u \times v = u_0 v_1 - u_1 v_0$,
  $u \cdot v = u_0 v_0 + u_1 v_1$, and
  $\lVert v\rVert^2 = v_0^2 + v_1^2$.
\end{definition}

\begin{lemma}[Squared length is nonnegative, and positive on nonzero vectors]
  \label{lem:lengthSq-pos}
  \lean{RationalPoint.lengthSq_nonneg, RationalPoint.lengthSq_pos_of_ne}
  \leanok
  \uses{def:rational-ops}
  For all $v \in \mathbb{Q}^2$, $\lVert v\rVert^2 \ge 0$, with strict
  inequality when $v \ne 0$.
\end{lemma}

\begin{definition}[Counterclockwise rotation by a right angle]
  \label{def:rotate90}
  \lean{RationalPoint.rotate90Counterclockwise}
  \leanok
  \uses{def:rational-point}
  The map $R : (x, y) \mapsto (-y, x)$ is a $90^\circ$ counterclockwise
  rotation, and preserves squared length.
\end{definition}

\begin{definition}[Closed / open half-spaces]
  \label{def:halfspace}
  \lean{ClosedHalfSpace, OpenHalfSpace}
  \leanok
  \uses{def:rational-point}
  A (closed, resp.\ open) \emph{half-space} is specified by a basepoint
  $b \in \mathbb{Q}^2$ and a nonzero inward normal $n \in \mathbb{Q}^2$;
  a point $p$ belongs to the half-space iff $n \cdot (p - b) \ge 0$
  (resp.\ $> 0$).
\end{definition}

\begin{definition}[Half-space to the left of a directed edge]
  \label{def:toLeft}
  \lean{RationalPoint.toStrictlyLeft, RationalPoint.toWeaklyLeft}
  \leanok
  \uses{def:halfspace, def:rotate90}
  Given distinct points $p_1, p_2 \in \mathbb{Q}^2$, the half-space
  strictly (resp.\ weakly) to the left of the directed segment
  $p_1 \to p_2$ uses basepoint $p_1$ and normal $R(p_2 - p_1)$.
\end{definition}

\begin{definition}[Lines and intersection]
  \label{def:line}
  \lean{Line, Line.intersection, ClosedHalfSpace.lineIntersection}
  \leanok
  \uses{def:rational-point, def:halfspace}
  A \emph{line} is a basepoint together with a nonzero direction vector.
  Two nonparallel lines have a unique intersection point, computed by
  Cramer's rule.
\end{definition}

\begin{definition}[Moving a half-space inward]
  \label{def:moveInward}
  \lean{ClosedHalfSpace.moveInward}
  \leanok
  \uses{def:halfspace, def:rational-utility}
  Given $d, \tau > 0$, one may translate the basepoint of a half-space
  inward along its normal by an amount in $[d, d+\tau]$, obtaining a new
  half-space contained in the original. The scaling factor is realised
  rationally via Lemma~\ref{lem:findRationalWithSquareBetween}.
\end{definition}

\begin{definition}[Rational root with square between bounds]
  \label{def:rational-utility}
  \lean{findRationalWithSquareBetween}
  \leanok
  Given $0 \le \ell < u$ in $\mathbb{Q}$, there is an explicit rational
  $r > 0$ with $\ell \le r^2 \le u$.
\end{definition}

\begin{lemma}[Correctness of \texttt{findRationalWithSquareBetween}]
  \label{lem:findRationalWithSquareBetween}
  \lean{findRationalWithSquareBetween_spec, findRationalWithSquareBetween_positive}
  \leanok
  \uses{def:rational-utility}
  The rational $r$ returned by \ref{def:rational-utility} satisfies $r > 0$,
  $\ell \le r^2$, and $r^2 \le u$.
\end{lemma}

\chapter{Convex polygons}

\begin{definition}[Non-degenerate polygon]
  \label{def:nondegen-polygon}
  \lean{NondegenPolygon}
  \leanok
  \uses{def:rational-point}
  A \emph{non-degenerate polygon} is a tuple consisting of a natural number
  $n \ge 3$, a function $V : \mathrm{Fin}\,n \to \mathbb{Q}^2$ giving the
  vertices in counterclockwise order, and a proof that $V$ is injective.
\end{definition}

\begin{definition}[Convex polygon]
  \label{def:convex-polygon}
  \lean{ConvexPolygon}
  \leanok
  \uses{def:nondegen-polygon, def:toLeft}
  A \emph{convex polygon} is a non-degenerate polygon in which, for every
  directed edge $V_i \to V_{i+1}$, every other vertex lies strictly to the
  left of the edge. Strictness excludes collinear triples and therefore
  makes every vertex an extreme point of the convex hull of $\{V_i\}$.
\end{definition}

\begin{definition}[Edges as half-spaces]
  \label{def:toHalfSpaces}
  \lean{ConvexPolygon.toHalfSpaces}
  \leanok
  \uses{def:convex-polygon, def:toLeft}
  Each edge of a convex polygon defines a closed half-space weakly to its
  left; the polygon coincides with the intersection of these half-spaces.
\end{definition}

\begin{definition}[Containment and subset tests]
  \label{def:contains}
  \lean{ConvexPolygon.contains, ConvexPolygon.isSubsetOf}
  \leanok
  \uses{def:toHalfSpaces}
  A point lies in a convex polygon iff it lies in every edge half-space.
  A convex polygon $P$ is a subset of $Q$ iff every vertex of $P$ lies in
  $Q$.
\end{definition}

\begin{definition}[Convex hull of a list of rational points]
  \label{def:convexHullRationalPoints}
  \lean{convexHullRationalPoints, ConvexPolygon.ofList}
  The function \texttt{convexHullRationalPoints} returns, from a finite
  list of rational points, a list of extreme points in counterclockwise
  order starting at the lexicographically smallest point.
  \texttt{ConvexPolygon.ofList} upgrades this to an \texttt{Option
  ConvexPolygon}, returning \texttt{none} when fewer than three extreme
  points remain.
\end{definition}

\begin{lemma}[Convex hull vertices are distinct]
  \label{lem:convexHullRationalPoints-nodup}
  \lean{convexHullRationalPoints_nodup}
  \uses{def:convexHullRationalPoints}
  The output of \texttt{convexHullRationalPoints} has no duplicates.
\end{lemma}

\begin{definition}[Intersection of half-spaces as a polygon]
  \label{def:ofHalfSpaces}
  \lean{ConvexPolygon.ofHalfSpaces}
  \leanok
  \uses{def:halfspace, def:convexHullRationalPoints, def:line}
  Given a list of closed half-spaces, the convex polygon defined by their
  intersection is obtained by taking all pairwise boundary intersections
  and retaining those satisfying every half-space constraint.
\end{definition}

\begin{definition}[Shrinking a convex polygon]
  \label{def:shrink}
  \lean{ConvexPolygon.shrink}
  \leanok
  \uses{def:toHalfSpaces, def:moveInward, def:ofHalfSpaces}
  Given $d, \tau > 0$, one shrinks a convex polygon by moving each edge
  inward by an amount in $[d, d + \tau]$ and intersecting the resulting
  half-spaces. The output is an \texttt{Option ConvexPolygon}, since the
  result may be empty.
\end{definition}

\begin{definition}[Shoelace area]
  \label{def:shoelace}
  \lean{shoelaceArea, ConvexPolygon.area}
  \leanok
  \uses{def:convex-polygon}
  The area of a convex polygon with vertex list $(V_0, \ldots, V_{n-1})$
  is
  \[
    \operatorname{area}(P) = \frac{1}{2}
    \left| \sum_{i=0}^{n-1} (V_i)_0 (V_{i+1})_1 - (V_{i+1})_0 (V_i)_1 \right|.
  \]
\end{definition}

\chapter{Planar isometries}

\begin{definition}[Direct planar isometry]
  \label{def:direct-isometry}
  \lean{Moser.DirectIsometry}
  \leanok
  \uses{def:rational-point}
  A \emph{direct isometry} is a tuple $(c, s, t)$ with $c, s \in \mathbb{Q}$,
  $t \in \mathbb{Q}^2$, and $c^2 + s^2 = 1$. It acts on $p \in \mathbb{Q}^2$
  by
  \[
    \mathrm{apply}_{(c,s,t)}(p) = (c\,p_0 - s\,p_1,\; s\,p_0 + c\,p_1) + t.
  \]
\end{definition}

\begin{lemma}[Direct isometries are bijections]
  \label{lem:isometry-bijective}
  \lean{Moser.DirectIsometry.apply_injective, Moser.DirectIsometry.apply_surjective}
  \leanok
  \uses{def:direct-isometry}
  The action of a direct isometry on $\mathbb{Q}^2$ is both injective and
  surjective.
\end{lemma}

\begin{definition}[Applying an isometry to a polygon]
  \label{def:applyPolygon}
  \lean{Moser.DirectIsometry.applyPolygon}
  \leanok
  \uses{def:direct-isometry, def:convex-polygon}
  Applying a direct isometry vertex-wise to a convex polygon produces a
  convex polygon. Convexity is preserved because rotations preserve the
  cross products witnessing the ``strictly left of'' condition.
\end{definition}

\begin{definition}[Composition of isometries]
  \label{def:isometry-compose}
  \lean{Moser.DirectIsometry.compose}
  \leanok
  \uses{def:direct-isometry}
  Direct isometries are closed under composition, with
  $(c, s, t)_1 \circ (c, s, t)_2$ given by the usual product of rotation
  matrices and composed translation.
\end{definition}

\begin{definition}[Rational unit circle grid]
  \label{def:angleGrid}
  \lean{Moser.angleGrid, Moser.angleGridAux}
  \leanok
  \uses{def:rational-point}
  For each pair of naturals $(a, b) \ne (0, 0)$, the formulas
  $c = (a^2 - b^2)/(a^2 + b^2)$ and $s = 2ab/(a^2 + b^2)$ yield a rational
  point on the unit circle. Enumerating such pairs up to a resolution
  depending on a target \emph{maximum angle change} yields a finite
  rational grid on the unit circle.
\end{definition}

\begin{lemma}[Angle grid lies on the unit circle]
  \label{lem:angleGrid-spec}
  \lean{Moser.angleGrid_spec, Moser.unit_circle_from_complex_square}
  \leanok
  \uses{def:angleGrid}
  Every $(c, s)$ produced by \texttt{angleGrid} satisfies $c^2 + s^2 = 1$.
\end{lemma}

\begin{definition}[Isometry discretisation]
  \label{def:discretizeIsometries}
  \lean{Moser.discretizeIsometries, Moser.defaultDiscretization}
  \leanok
  \uses{def:direct-isometry, def:angleGrid, def:locationRange}
  For a granularity $\epsilon > 0$, we form a finite list of direct
  isometries by pairing each angle from \texttt{angleGrid} with each
  translation on a rational grid intersected with the
  \texttt{LocationRange} (Definition~\ref{def:locationRange}). The angle
  resolution is scaled by the distance cutoff so that far-away points are
  still covered finely.
\end{definition}

\chapter{Constants and the search region}

\begin{definition}[Area threshold]
  \label{def:areaThreshold}
  \lean{Moser.areaThreshold}
  \leanok
  The area threshold is $A_* := 232240 / 10^6 = 0.232240$. A Moser set
  with area strictly less than $A_*$ would improve the current best lower
  bound.
\end{definition}

\begin{definition}[Benchmark worms]
  \label{def:benchmark-worms}
  \lean{Moser.IsocelesRightTriangleWorm, Moser.SquareWorm, Moser.RightTriangleOneThirdWorm, Moser.InitialWorm}
  \leanok
  \uses{def:convex-polygon}
  Three concrete convex polygons are fixed as benchmark worms:
  \begin{itemize}
    \item the isosceles right triangle with legs of length $1/2$,
    \item the unit square with side $1/3$,
    \item the right triangle with legs $1/3$ and $2/3$.
  \end{itemize}
  The \emph{initial worm} is taken to be the isosceles right triangle; by
  convention every candidate Moser set in the working set contains an
  un-shifted copy of it. Its area is exactly $1/8$.
\end{definition}

\begin{definition}[Location range and distance cutoff]
  \label{def:locationRange}
  \lean{Moser.offset, Moser.LocationRange, Moser.upperBoundSqrtTwo, Moser.distanceCutoff}
  \leanok
  \uses{def:areaThreshold, def:convex-polygon}
  Let $\sigma := 4 A_*$. The \emph{location range} is the square
  $[-\sigma, \sigma]^2$, which upper-bounds the positions where a point of
  a working-set polygon may lie without already exceeding the area
  threshold (via a triangle formed with the initial worm). The
  \emph{distance cutoff} is $\sigma \cdot \overline{\sqrt{2}}$, where
  $\overline{\sqrt{2}} = 1414213562373095/10^{15}$ is a rational upper
  bound on $\sqrt 2$.
\end{definition}

\begin{proposition}[Points outside the location range are excluded]
  \label{prop:locationRange-excludes}
  \lean{Moser.area_hull_initialWorm_insert_gt_areaThreshold}
  \uses{def:locationRange, def:areaThreshold, def:benchmark-worms}
  Any convex polygon containing the initial worm as well as a point
  outside the location range has area strictly greater than $A_*$.
  Concretely, if a rational point $p$ is not contained in the location
  range, then the convex hull of $\{p\}$ together with the vertices of
  the initial worm has area strictly greater than $A_*$.
\end{proposition}
\begin{proof}
  \uses{def:locationRange, def:areaThreshold, def:benchmark-worms}
  By definition the location range is the square with corners at
  $(\pm 4 A_*, \pm 4 A_*)$. If $p = (x, y)$ lies outside this square,
  then $\max(|x|, |y|) > 4 A_*$. The initial worm has vertices
  $(0,0)$, $(1/2, 0)$, $(0, 1/2)$, so the triangle formed by $p$
  together with one of the legs of the initial worm (the leg of
  length $1/2$ along the axis where $p$ is far from the origin) has
  base $1/2$ and height greater than $4 A_*$, hence area greater than
  $A_*$. This triangle is contained in the convex hull of $p$ together
  with the worm vertices, so the hull's area exceeds $A_*$.
\end{proof}

\chapter{Worms}

\begin{definition}[Newton-style rational square root]
  \label{def:sqrtApprox}
  \lean{Moser.sqrtApprox, Moser.distanceApprox}
  \leanok
  \uses{def:rational-point}
  Given $s \ge 0$ and $\epsilon > 0$, the Newton iteration
  $x_{n+1} = (x_n + s/x_n)/2$ produces a rational $r$ with
  $|r - \sqrt{s}| < \epsilon$. Applied componentwise, this yields a
  rational approximation \texttt{distanceApprox} of the Euclidean distance
  between two rational points.
\end{definition}

\begin{definition}[Worm]
  \label{def:worm}
  \lean{Moser.Worm}
  \leanok
  \uses{def:rational-point}
  A \emph{worm} is a polygonal path with at least two rational vertices.
  Its approximate total length is the sum of distance approximations along
  consecutive vertices, with per-segment error budget divided equally so
  the total error remains below a given $\epsilon$.
\end{definition}

\begin{definition}[Scaling a worm, unit worm]
  \label{def:scaleToUnit}
  \lean{Moser.Worm.scale, Moser.Worm.scaleToUnit, Moser.UnitWorm, Moser.Worm.toUnitWorm}
  \uses{def:worm, def:sqrtApprox}
  A worm can be scaled by any rational factor; scaling by $1/\ell$ where
  $\ell$ is the approximate length brings the worm to approximately unit
  length. A \emph{unit worm} bundles a worm together with a uniform bound
  ensuring the approximate length tends to $1$ as $\epsilon \to 0$.
\end{definition}

\begin{definition}[Convex hull of a worm]
  \label{def:worm-hull}
  \lean{Moser.Worm.toConvexPolygon, Moser.Worm.convexHullPolygon, Moser.UnitWorm.toConvexPolygon}
  \uses{def:worm, def:convexHullRationalPoints}
  The convex hull of a worm is a convex polygon (when the worm has at
  least three non-collinear vertices); it summarises all rigid motions of
  the worm that embed it into a convex target.
\end{definition}

\chapter{Moser sets}

\begin{definition}[Moser set]
  \label{def:moser-set}
  A set $M \subseteq \mathbb{R}^2$ is a \emph{Moser set} if for every
  worm $w$ of unit length there exists a direct isometry $\varphi$ with
  $\varphi(w) \subseteq M$.
\end{definition}

\begin{definition}[Approximate / computational Moser set]
  \label{def:approx-moser-set}
  \uses{def:worm, def:direct-isometry}
  Given a finite list of candidate worms $W$ and a finite list of
  isometries $I$, a set $M$ is a \emph{$(W, I)$-approximate Moser set} if
  for every $w \in W$ there exists $\varphi \in I$ with
  $\varphi(w) \subseteq M$.
\end{definition}

\chapter{Working set algorithm}

The computational approach maintains a list of convex polygons, the
\emph{working set}, with the invariant that every Moser set of area
below the threshold contains (a rigid motion of) at least one polygon in
the list.

\begin{definition}[Working set]
  \label{def:workingSet}
  \lean{Moser.WorkingSet}
  \leanok
  \uses{def:convex-polygon}
  A \emph{working set} is a list of convex polygons. Intended invariants:
  \begin{enumerate}
    \item every polygon in the list contains a copy of the initial worm
      (Definition~\ref{def:benchmark-worms}) under some direct isometry;
    \item every Moser set of area strictly less than $A_*$ contains a
      copy of some polygon in the list under some direct isometry.
  \end{enumerate}
  Invariant 1 is currently a convention maintained by construction;
  Invariant 2 is the core correctness statement the algorithm preserves.
\end{definition}

\begin{definition}[Initial working set]
  \label{def:initialWorkingSet}
  \lean{Moser.WorkingSet.initial, Moser.WorkingSet.InitialWorkingSet}
  \leanok
  \uses{def:workingSet, def:benchmark-worms}
  The initial working set contains exactly the initial worm.
\end{definition}

\begin{definition}[Minimum area]
  \label{def:minArea}
  \lean{Moser.WorkingSet.minAreaPolygon, Moser.WorkingSet.minArea, Moser.WorkingSet.size}
  \leanok
  \uses{def:workingSet, def:shoelace}
  The minimum area in a working set is the minimum of the shoelace areas
  of its polygons (taken as $0$ on the empty set).
\end{definition}

\chapter{Convex hull growth}

Given an area threshold $A_*$ and a nontrivial convex polygon of area below
the threshold, there is a limited region of the plane such that adjoining a
point in that region to the polygon and taking the convex hull keeps the
resulting area below $A_*$. Our goal in this chapter is to describe this
region as the intersection of a finite collection of half-spaces, one for
each ordered pair of vertices of the input polygon, computed up to a
specified rational tolerance $\tau$.

\begin{definition}[Growable distance to the left of a vertex pair]
  \label{def:growableDistance}
  \lean{ConvexPolygon.growableDistance}
  \uses{def:convex-polygon, def:areaThreshold}
  Let $P$ be a convex polygon, $A_*$ a rational area threshold, and
  $V_i, V_j$ an ordered pair of (not necessarily adjacent) vertices of
  $P$. Let $\tau > 0$ be a rational tolerance. The
  \emph{growable distance to the left of the directed line
  $V_i \to V_j$} is a rational number $d \ge 0$ such that, writing $A_R$
  for the area of $P$ lying weakly to the right of the directed line
  $V_i V_j$ and $T(d)$ for the area of the triangle formed by $V_i$,
  $V_j$, and a point at perpendicular distance $d$ to the left of the
  line, we have
  \[ A_R + T(d) \;\ge\; A_* , \]
  and $d$ is within the tolerance $\tau$ of the unique exact value for
  which equality holds.
\end{definition}

\begin{definition}[Growth half-space to the left of a vertex pair]
  \label{def:growthHalfspace}
  \lean{ConvexPolygon.growthHalfspace, ConvexPolygon.growthHalfspaceOfPair}
  \uses{def:growableDistance, def:halfspace}
  Let $P$, $A_*$, $\tau$, and $V_i, V_j$ be as in
  Definition~\ref{def:growableDistance}, with growable distance $d$.
  The \emph{growth half-space} to the left of the pair $V_i, V_j$ is
  the closed half-space lying to the left of (and including) the line
  parallel to $V_i \to V_j$ at perpendicular distance $d$ on the left
  side of the directed line.
\end{definition}

\begin{lemma}[Threshold violated outside a growth half-space]
  \label{lem:threshold_violated_outside_growth_halfspace}
  \lean{ConvexPolygon.threshold_violated_outside_growth_halfspace}
  \uses{def:growthHalfspace}
  Let $P$ be a convex polygon with rational vertices, $A_*$ a rational
  area threshold, $\tau > 0$ a rational tolerance, and $V_i, V_j$ an
  ordered pair of vertices of $P$. If a rational point $p$ lies
  strictly outside the growth half-space to the left of $V_i, V_j$
  associated with $P$, $A_*$, and $\tau$, then the area of the convex
  hull of $P \cup \{p\}$ is strictly greater than $A_*$.
\end{lemma}
\begin{proof}
  \uses{def:growableDistance, def:growthHalfspace}
  By the definition of the growth half-space, $p$ lying strictly
  outside means that the perpendicular signed distance from $p$ to the
  directed line $V_i V_j$, measured to the left, exceeds the growable
  distance $d$. The convex hull of $P$ together with $p$ contains both
  the portion of $P$ lying weakly to the right of the line $V_i V_j$
  (of area $A_R$) and the triangle with vertices $V_i$, $V_j$, $p$,
  and these two regions intersect only along the segment $V_i V_j$, a
  set of measure zero. The triangle has base $|V_i V_j|$ and height
  equal to the left-distance from $p$ to the line, hence area strictly
  greater than $T(d)$. Therefore the hull's area is at least
  $A_R + T(\text{left-distance of } p) > A_R + T(d) \ge A_*$.
\end{proof}

\begin{definition}[Growth half-space intersection polygon]
  \label{def:growthHalfspaceIntersection}
  \lean{ConvexPolygon.growthHalfspaceIntersection}
  \uses{def:growthHalfspace, def:ofHalfSpaces}
  Let $P$ be a convex polygon, $A_*$ a rational area threshold, and
  $\tau > 0$ a rational tolerance. The
  \emph{growth half-space intersection} associated to $P$, $A_*$, and
  $\tau$ is the intersection, over all ordered pairs $(V_i, V_j)$ of
  distinct vertices of $P$, of the growth half-spaces to the left of
  those pairs. Equivalently, it is the set of points $p$ such that for
  every ordered vertex pair $(V_i, V_j)$, the point $p$ lies in the
  growth half-space to the left of $V_i, V_j$.
\end{definition}

\begin{lemma}[Threshold violated outside the growth half-space intersection]
  \label{lem:threshold_violated_outside_growth_intersection}
  \lean{ConvexPolygon.threshold_violated_outside_growth_intersection}
  \uses{def:growthHalfspaceIntersection}
  Let $P$ be a convex polygon with rational vertices, $A_*$ a rational
  area threshold, and $\tau > 0$ a rational tolerance. If a rational
  point $p$ lies outside the growth half-space intersection of $P$,
  $A_*$, and $\tau$, then the area of the convex hull of
  $P \cup \{p\}$ is strictly greater than $A_*$.
\end{lemma}
\begin{proof}
  \uses{lem:threshold_violated_outside_growth_halfspace,
    def:growthHalfspaceIntersection}
  If $p$ is outside the intersection, then there exists at least one
  ordered vertex pair $(V_i, V_j)$ such that $p$ lies strictly outside
  the growth half-space to the left of $(V_i, V_j)$. By
  Lemma~\ref{lem:threshold_violated_outside_growth_halfspace}, the
  convex hull of $P \cup \{p\}$ then has area strictly greater than
  $A_*$.
\end{proof}



\section{Pruning operations}

\begin{definition}[Big set removal]
  \label{def:bigSetRemoval}
  \lean{Moser.WorkingSet.bigSetRemoval}
  \leanok
  \uses{def:workingSet, def:areaThreshold, def:shoelace}
  Remove every polygon whose area strictly exceeds $A_*$. This preserves
  Invariant~2 because any Moser set of area below $A_*$ cannot have a
  polygon of strictly larger area as a subset.
\end{definition}

\begin{proposition}[Big set removal preserves the working set invariant]
  \label{prop:bigSetRemoval-correct}
  \uses{def:bigSetRemoval, def:moser-set}
  If a working set satisfies Invariant~2 of
  Definition~\ref{def:workingSet}, so does the result of big set removal.
\end{proposition}

\begin{definition}[Superset removal]
  \label{def:supersetRemoval}
  \lean{Moser.WorkingSet.supersetRemoval}
  \leanok
  \uses{def:workingSet, def:contains}
  Remove every polygon $p$ in the list for which some distinct polygon
  $q$ in the list is a subset of $p$. The smaller $q$ already witnesses
  Invariant~2 whenever $p$ does.
\end{definition}

\begin{proposition}[Superset removal preserves the working set invariant]
  \label{prop:supersetRemoval-correct}
  \uses{def:supersetRemoval, def:moser-set}
  If a working set satisfies Invariant~2, so does the result of superset
  removal.
\end{proposition}

\begin{definition}[Worm adding]
  \label{def:wormAdding}
  \lean{Moser.WorkingSet.wormAdding}
  \leanok
  \uses{def:workingSet, def:discretizeIsometries, def:applyPolygon, def:shrink, def:convexHullRationalPoints}
  Given an additional worm presented as a convex hull, for each polygon
  $p$ in the current working set and for each isometry $\varphi$ in the
  discretisation, form the convex hull of $p$ together with the image
  under $\varphi$ of the \emph{shrunk} hull (to absorb discretisation
  error). The new working set collects all such unions. This step encodes
  the requirement that a Moser set must contain \emph{this} worm in some
  pose.
\end{definition}

\begin{proposition}[Worm adding preserves the working set invariant]
  \label{prop:wormAdding-correct}
  \uses{def:wormAdding, def:moser-set, def:discretizeIsometries}
  If the isometry discretisation is fine enough relative to the shrink
  amount, then adding a worm preserves Invariant~2 of
  Definition~\ref{def:workingSet}.
\end{proposition}

\begin{definition}[Cleanup and combined add]
  \label{def:cleanup}
  \lean{Moser.WorkingSet.cleanup, Moser.WorkingSet.addWormAndCleanup}
  \leanok
  \uses{def:bigSetRemoval, def:supersetRemoval, def:wormAdding}
  \texttt{cleanup} is the composition of big-set removal and superset
  removal. \texttt{addWormAndCleanup} adds a worm and then cleans up.
\end{definition}




\chapter{Lower bound theorem}

The end goal of the working set iteration is the following claim.

\begin{theorem}[Computational lower bound]
  \label{thm:lower-bound}
  \uses{def:moser-set, def:areaThreshold, def:workingSet, def:cleanup}
  There exists a finite list of worms $W$ and a sequence of working-set
  updates that, starting from
  Definition~\ref{def:initialWorkingSet} and adding the worms in $W$ in
  turn using
  Definition~\ref{def:cleanup}, terminates with the empty working set.
  Consequently the area of every Moser set is at least $A_*$:
  \[
    \inf \{ \operatorname{area}(M) : M \text{ is a Moser set}\} \ge A_*.
  \]
\end{theorem}

\begin{proof}
  \uses{prop:bigSetRemoval-correct, prop:supersetRemoval-correct, prop:wormAdding-correct, prop:locationRange-excludes}
  Any Moser set $M$ of area $<A_*$ must contain, in some pose, every
  polygon ever present in a working set produced by
  \texttt{addWormAndCleanup}, by Propositions
  \ref{prop:bigSetRemoval-correct},
  \ref{prop:supersetRemoval-correct}, and \ref{prop:wormAdding-correct}.
  If the final working set is empty, this is a contradiction.
\end{proof}

\begin{proposition}[Terminating sequence of worms]
  \label{prop:terminating-sequence}
  \uses{thm:lower-bound, def:benchmark-worms}
  A finite list of worms ending the iteration with the empty working
  set can be found by computer search, e.g.\ starting from the
  benchmark worms of Definition~\ref{def:benchmark-worms} and augmenting
  with further worms arising from the enumeration scheme in
  \texttt{Moser.Worm.Enumeration} (currently stubbed).
\end{proposition}

\chapter{Compactness and the Moser cover number}

This chapter develops an alternative, real-analytic framing of the problem
that does not rely on the rational discretisation of the preceding chapters.
A worm is taken to be an arbitrary $1$-Lipschitz curve, and we approximate
the (infinite) space of worms by a finite $\varepsilon$-net of grid
polygonal worms. This yields matching upper and lower bounds on the
\emph{Moser cover number}, whose gap is controlled by the planar Steiner
formula. The development is carried out in
\texttt{Moser/Real/CompactnessOutline.lean}.

\section{Worms and thickenings}

\begin{definition}[Worm (Lipschitz curve)]
    \label{def:lipschitz-worm}
    \lean{Moser.CompactnessOutline.Worms}
    A \emph{worm} is a $1$-Lipschitz curve $f : [0,1] \to \RR^2$.
\end{definition}

\begin{definition}[Pinned worm]
    \label{def:pinnedWorm}
    \lean{Moser.CompactnessOutline.PinnedWorms}
    \uses{def:lipschitz-worm}
    A worm $f$ is \emph{pinned} if $f(0) = (0,0)$. Since $f$ is $1$-Lipschitz
    on $[0,1]$, a pinned worm satisfies $\lvert f(x) \rvert \le x \le 1$ for
    all $x \in [0,1]$.
\end{definition}

\begin{definition}[$\varepsilon$-thickening]
    \label{def:thickening}
    \lean{Metric.cthickening}
    For a set $A \subseteq \RR^2$ and $\varepsilon \ge 0$, the
    \emph{$\varepsilon$-thickening} of $A$ is the closed
    $\varepsilon$-neighbourhood
    \[
        A^{\varepsilon} \;=\; \{\, x \in \RR^2 : \dist(x, A) \le \varepsilon \,\}.
    \]
\end{definition}

\begin{definition}[Grid polygonal worm]
    \label{def:gridPolygonalWorm}
    \lean{Moser.CompactnessOutline.GridPolygonalWorms}
    \uses{def:pinnedWorm}
    Fix a step $\delta > 0$ and a vertex count $k \in \NN$. A
    \emph{$(k,\delta)$-grid polygonal worm} is a pinned worm obtained by
    choosing grid points $p_0, p_1, \dots, p_k \in \delta\ZZ \times \delta\ZZ$
    with $p_0 = (0,0)$, setting $\tilde f(n/k) = p_n$ for $0 \le n \le k$, and
    interpolating linearly on each interval $[n/k, (n+1)/k]$, subject to the
    resulting curve being $1$-Lipschitz.
\end{definition}

\section{A finite \texorpdfstring{$\varepsilon$}{epsilon}-net of polygonal worms}

\begin{definition}[$\varepsilon$-net of worms]
    \label{def:epsilonNet}
    \lean{Moser.CompactnessOutline.IsEpsilonNet}
    \uses{def:pinnedWorm, def:thickening}
    Let $\varepsilon > 0$. A set $S$ of worms is an \emph{$\varepsilon$-net of
    worms} if the image of every pinned worm is contained in the
    $\varepsilon$-thickening of the image of some element of $S$.
\end{definition}

\begin{theorem}[A finite grid $\varepsilon$-net exists]
    \label{thm:finiteEpsilonNet}
    \lean{Moser.CompactnessOutline.finiteEpsilonNet}
    \uses{def:gridPolygonalWorm, def:epsilonNet}
    For every $\varepsilon > 0$ there exist $k \in \NN$ and $\delta > 0$ such
    that the set $S_\varepsilon$ of $(k,\delta)$-grid polygonal worms whose
    nodes lie within distance $1$ of the origin is finite and is an
    $\varepsilon$-net of worms.
\end{theorem}
\begin{proof}
    \uses{def:pinnedWorm, def:gridPolygonalWorm, def:thickening, def:epsilonNet}
    Choose $k > 2/\varepsilon$, so that $1/k < \varepsilon/2$, and then choose
    $\delta > 0$ small enough that $2\delta k + \delta < \varepsilon/2$. With
    these $k$ and $\delta$, the set $S_\varepsilon$ is finite, since each of the
    finitely many nodes ranges over the finitely many grid points of
    $\delta\ZZ \times \delta\ZZ$ inside the unit disc.

    It remains to check that $S_\varepsilon$ is an $\varepsilon$-net. Fix a
    pinned worm $f$ and construct an element of $S_\varepsilon$ near it. Set
    \[
        f_1 = \frac{1}{1 + 2\delta k}\, f ,
    \]
    which is $\tfrac{1}{1+2\delta k}$-Lipschitz; since $f$ is pinned,
    $\lvert f(x) \rvert \le 1$, so
    \[
        \lvert f(x) - f_1(x) \rvert
            = \lvert f(x) \rvert \cdot \frac{2\delta k}{1 + 2\delta k}
            \le 2\delta k .
    \]
    Let $\tilde f(n/k)$ be a nearest point of $\delta\ZZ \times \delta\ZZ$ to
    $f_1(n/k)$ for $0 \le n \le k$, and extend $\tilde f$ by linear
    interpolation. Each node moves by at most $\delta$ (half a grid cell in
    each coordinate), so on each subinterval of length $1/k$ the segment of
    $\tilde f$ has length at most
    $\tfrac{1}{k}\cdot\tfrac{1}{1+2\delta k} + 2\delta$; multiplying by $k$
    shows the Lipschitz constant of $\tilde f$ is at most
    $\tfrac{1}{1+2\delta k}\,(1 + 2\delta k) = 1$, so $\tilde f$ is a
    $(k,\delta)$-grid polygonal worm, and its nodes lie within distance $1$ of
    the origin, hence $\tilde f \in S_\varepsilon$. Comparing $f$ to $\tilde f$
    at an arbitrary $x \in [0,1]$, the rescaling contributes at most
    $2\delta k$, the grid rounding at most $\delta$, and the linear
    interpolation at most $1/k$, so
    \[
        \lvert f(x) - \tilde f(x) \rvert
            \;\le\; 2\delta k + \delta + \frac{1}{k} < \varepsilon .
    \]
    Therefore every point of the image of $f$ lies within $\varepsilon$ of the
    image of $\tilde f$, i.e.\ the image of $f$ is contained in the
    $\varepsilon$-thickening of the image of $\tilde f \in S_\varepsilon$.
\end{proof}

\section{Bounds on the Moser number}

\begin{definition}[Moser cover number]
    \label{def:moserNumber}
    \lean{Moser.CompactnessOutline.moserCoverNumber}
    \uses{def:lipschitz-worm}
    A set $C \subseteq \RR^2$ is a \emph{cover} if for every worm $f$ there is
    an orientation-preserving isometry (a translation--rotation) $g$ of $\RR^2$
    such that the image of $g \circ f$ is contained in $C$. The \emph{Moser
    cover number} $M$ is the infimum of $\area(C)$ over all convex covers $C$.
\end{definition}

\begin{definition}[Minimal convex cover of a finite set of worms]
    \label{def:minimalCover}
    \lean{Moser.CompactnessOutline.IsMinimalCover}
    \uses{def:lipschitz-worm}
    For a finite set $S$ of worms, the \emph{minimal convex cover} $K(S)$ is a
    convex set of least area for which there exist translation--rotations
    $(g_s)_{s \in S}$ with the image of $g_s \circ s$ contained in $K(S)$ for
    every $s \in S$; that is, $K(S)$ minimises
    $\area\!\big(\conv \bigcup_{s \in S} g_s(\image\, s)\big)$ over all choices
    of translation--rotations $g_s$.
\end{definition}

\begin{remark}
    \label{rem:computeMinimalCover}
    \uses{def:minimalCover}
    The minimal convex cover $K(S)$ is the minimum, over all
    translation--rotations of the elements of $S$, of the area of the convex
    hull of their union. For a fixed finite $S$ this is a finite-dimensional
    optimisation problem in the placement parameters and is in principle
    computable, e.g.\ by quantifier elimination over the reals, and plausibly
    by convex optimisation.
\end{remark}

\begin{theorem}[Lower bound on the Moser number]
    \label{thm:lowerBound}
    \lean{Moser.CompactnessOutline.lowerBound}
    \uses{def:moserNumber, def:minimalCover}
    For any finite set $S$ of worms, $M \ge \area\big(K(S)\big)$.
\end{theorem}
\begin{proof}
    \uses{def:moserNumber, def:minimalCover}
    Let $C$ be any convex cover. Since every element of $S$ is a worm, for each
    $s \in S$ there is a translation--rotation $g_s$ placing $g_s \circ s$
    inside $C$. Thus $C$ is a convex set admitting placements of all elements of
    $S$, so by minimality $\area(C) \ge \area(K(S))$. Taking the infimum over
    convex covers $C$ gives $M \ge \area(K(S))$.
\end{proof}

\begin{theorem}[Upper bound on the Moser number]
    \label{thm:upperBound}
    \lean{Moser.CompactnessOutline.upperBound}
    \uses{def:moserNumber, def:minimalCover, def:thickening, def:epsilonNet}
    Let $\varepsilon > 0$ and let $S_\varepsilon$ be a finite
    $\varepsilon$-net of worms. Then
    \[
        M \le \area\!\big(K(S_\varepsilon)^{\varepsilon}\big),
    \]
    the area of the $\varepsilon$-thickening of the minimal convex cover of
    $S_\varepsilon$.
\end{theorem}
\begin{proof}
    \uses{def:moserNumber, def:minimalCover, def:thickening, thm:finiteEpsilonNet}
    Write $K = K(S_\varepsilon)$ with placements $(g_s)_{s \in S_\varepsilon}$.
    The $\varepsilon$-thickening $K^{\varepsilon}$ is convex. We claim it is a
    cover. Let $f$ be an arbitrary worm; after a translation--rotation we may
    assume $f$ is pinned. By Theorem~\ref{thm:finiteEpsilonNet} there is
    $s \in S_\varepsilon$ such that the image of $f$ is contained in the
    $\varepsilon$-thickening of the image of $s$. Applying the placement $g_s$,
    the image of $g_s \circ f$ is contained in the $\varepsilon$-thickening of
    the image of $g_s \circ s$, which lies in $K^{\varepsilon}$. Hence
    $K^{\varepsilon}$ is a convex cover, and $M \le \area(K^{\varepsilon})$.
\end{proof}

\begin{theorem}[Steiner gap between the bounds]
    \label{thm:steinerGap}
    \lean{Moser.CompactnessOutline.steinerGap}
    \uses{def:minimalCover, def:thickening}
    Let $K$ be a convex body in $\RR^2$ with perimeter $L$. Then for every
    $\varepsilon \ge 0$,
    \[
        \area\big(K^{\varepsilon}\big)
            = \area(K) + \varepsilon\, L + \pi \varepsilon^{2}.
    \]
    In particular the upper and lower bounds of
    Theorems~\ref{thm:lowerBound} and~\ref{thm:upperBound} differ by at most
    $\varepsilon\, L + \pi \varepsilon^{2}$, where $L$ is the perimeter of
    $K(S_\varepsilon)$.
\end{theorem}
\begin{proof}
    \uses{def:thickening}
    This is the planar Steiner formula for convex bodies: the outer parallel
    body at distance $\varepsilon$ adds a strip of area $\varepsilon L$ along
    the boundary together with disc sectors of total area $\pi \varepsilon^2$ at
    the convex corners.
\end{proof}

\begin{remark}[Perimeter bound]
    \label{rem:perimeterBound}
    \uses{def:minimalCover, thm:steinerGap}
    To turn Theorem~\ref{thm:steinerGap} into an effective error bound one needs
    an a priori bound on the perimeter $L$ of $K(S_\varepsilon)$. Insisting that
    $S_\varepsilon$ contains a fixed closed worm (a loop, e.g.\ a small circle)
    forces every cover to be bounded, hence $K(S_\varepsilon)$ is bounded with a
    controlled perimeter. A quantitative perimeter bound is left as a
    \textsc{todo}; one may also add a square to $S_\varepsilon$ to make this
    argument go through.
\end{remark}

\begin{theorem}[Approximating the Moser number]
    \label{thm:approxAlgorithm}
    \lean{Moser.CompactnessOutline.approxAlgorithm}
    \uses{def:moserNumber, thm:finiteEpsilonNet, thm:lowerBound, thm:upperBound, thm:steinerGap, rem:perimeterBound}
    Assume a perimeter bound $L_0$ on $K(S_\varepsilon)$ as in
    Remark~\ref{rem:perimeterBound}. Then for any target accuracy $x > 0$ the
    Moser cover number $M$ can be computed to within $x$ by the following
    procedure:
    \begin{enumerate}
        \item choose $\varepsilon = x$ small enough that
            $\varepsilon L_0 + \pi \varepsilon^2 \le x$;
        \item compute a finite $\varepsilon$-net $S_\varepsilon$, for instance
            by repeatedly locating worms not yet covered and adding them;
        \item compute the minimal convex cover $K(S_\varepsilon)$.
    \end{enumerate}
    Then $\area(K(S_\varepsilon))$ is a lower bound and
    $\area(K(S_\varepsilon)^{\varepsilon})$ an upper bound for $M$, and the two
    differ by at most $x$.
\end{theorem}
\begin{proof}
    \uses{thm:lowerBound, thm:upperBound, thm:steinerGap}
    By Theorem~\ref{thm:lowerBound}, $\area(K(S_\varepsilon)) \le M$, and by
    Theorem~\ref{thm:upperBound}, $M \le \area(K(S_\varepsilon)^{\varepsilon})$.
    By Theorem~\ref{thm:steinerGap} the gap between these bounds equals
    $\varepsilon L + \pi \varepsilon^2 \le \varepsilon L_0 + \pi \varepsilon^2
    \le x$. Hence either bound approximates $M$ to within $x$.
\end{proof}

